\section{Evaluation Results and Insights}
\label{sec:results}

Square-bracket numbers in this section are placeholders. They are replaced
by the fields of \path{source_code/outputs/metrics.json}. The pipeline
does not tune on \(T_{\mathrm{test}}\).

\subsection{Model Performance Evaluation}
\label{sec:performance}

On a test block of \(n\) hours with demand \(y_t\) and forecast
\(\hat{y}_t\), the three error scores are
\begin{align}
  \mathrm{MAE} &= \frac{1}{n}\sum_{t=1}^{n}\lvert y_t-\hat{y}_t\rvert, \label{eq:mae}\\
  \mathrm{RMSE} &= \sqrt{\frac{1}{n}\sum_{t=1}^{n}(y_t-\hat{y}_t)^2}, \label{eq:rmse}\\
  \mathrm{WAPE} &= \frac{\sum_{t=1}^{n}\lvert y_t-\hat{y}_t\rvert}{\sum_{t=1}^{n}\lvert y_t\rvert}. \label{eq:wape}
\end{align}
MAE and RMSE are reported in kW. RMSE penalises large misses more than MAE,
which matters for hours near the capacity limit. WAPE scales the absolute
error by total demand, so it can be read as a system-wide percentage.

\begin{table}[h]
  \centering
  \small
  \begin{tabular}{@{}lrrr@{}}
    \toprule
    Model & MAE (kW) & RMSE (kW) & WAPE (\%) \\
    \midrule
    Seasonal naive & [XX] & [XX] & [XX] \\
    Ridge & [XX] & [XX] & [XX] \\
    Random forest & [XX] & [XX] & [XX] \\
    XGBoost & [XX] & [XX] & [XX] \\
    LightGBM & [XX] & [XX] & [XX] \\
    \bottomrule
  \end{tabular}
  \caption{Test-block scores. Source: \texttt{test\_scores} in \texttt{metrics.json}.}
  \label{tab:scores}
\end{table}

\begin{figure}[h]
  \centering
  \figplaceholder{Bar chart of test MAE and RMSE (\texttt{model\_comparison.png})\\
  and a two-week overlay of actual demand, the best forecast, and \(\tau\)\\
  (\texttt{forecast\_test.png}).}
  \caption{Error comparison and a visual check on the opening of \(T_{\mathrm{test}}\).}
  \label{fig:comparison}
\end{figure}

The improvement of a learned model against the seasonal naive is
\[
  \frac{\mathrm{MAE}_{\mathrm{naive}}-\mathrm{MAE}_{\mathrm{model}}}{\mathrm{MAE}_{\mathrm{naive}}}.
\]
Report that percentage from \texttt{mae\_improvement\_vs\_naive\_pct}. The
discussion should state whether LightGBM or XGBoost reduces both typical
error (MAE, WAPE) and large error (RMSE) relative to the weekly baseline.
A lower MAE alone is not enough if the peak hours remain biased.

\subsection{Peak Capacity and Risk Metrics}
\label{sec:peak}

Average error can look acceptable while the forecast still misses the hour
that threatens the network. Three peak scores are therefore computed on
\(T_{\mathrm{test}}\) only.

\paragraph{Peak magnitude error.}
For each test day, take the actual peak and the forecast peak and record the
absolute gap in kW. The reported value is the mean gap over days that contain
at least 12 hours.

\paragraph{Peak timing error.}
For the same days, record the circular distance in hours between the clock
time of the actual peak and the clock time of the forecast peak. A timing
error of one hour still leaves a short window in which to arm reserve. An
error of several hours does not.

\paragraph{Breach anticipation rate.}
Let \(\tau\) be the training 95th percentile, stored as \texttt{threshold\_kw}
(numeric value [XX] kW). A breach hour is an hour with \(y_t>\tau\). The
breach anticipation rate is the share of those hours for which the forecast
also exceeds \(\tau\):
\begin{equation}
  \mathrm{BAR}
  = \frac{\lvert\{t\in T_{\mathrm{test}}: y_t>\tau \ \mathrm{and}\ \hat{y}_t>\tau\}\rvert}
         {\lvert\{t\in T_{\mathrm{test}}: y_t>\tau\}\rvert}.
  \label{eq:bar}
\end{equation}
The one-hour-ahead forecast is issued before the hour closes, so a hit in
\eqref{eq:bar} is an alert that precedes the realised breach. Alert precision,
the share of threshold crossings that were real, is stored beside BAR so a
high hit rate is not confused with a model that alerts every hour.

\begin{table}[h]
  \centering
  \small
  \begin{tabular}{@{}lrrrr@{}}
    \toprule
    Model & Peak magnitude (kW) & Peak timing (h) & BAR (\%) & Alert precision (\%) \\
    \midrule
    Seasonal naive & [XX] & [XX] & [XX] & [XX] \\
    Ridge & [XX] & [XX] & [XX] & [XX] \\
    Random forest & [XX] & [XX] & [XX] & [XX] \\
    XGBoost & [XX] & [XX] & [XX] & [XX] \\
    LightGBM & [XX] & [XX] & [XX] & [XX] \\
    \bottomrule
  \end{tabular}
  \caption{Peak scores on \(T_{\mathrm{test}}\). Source: \texttt{peak} in \texttt{metrics.json}.}
  \label{tab:peak}
\end{table}

\subsection{Findings and Operating Implications}
\label{sec:insights}

Two exogenous drivers are the focus of the interpretation. Hour of day sets
the shape of the daily cycle. Temperature, together with the cooling-degree
and squared-temperature terms, sets how high that cycle rises on a hot day.
Lag and rolling features will often take a large share of LightGBM gain,
because a one-hour-ahead forecast is strongly autoregressive. The importance
figure should therefore be read in two layers: overall gain, and gain inside
the weather and calendar group. Fill the ranking from
\texttt{lightgbm\_importance} and \texttt{lightgbm\_importance\_groups}.

\begin{figure}[h]
  \centering
  \figplaceholder{Top LightGBM features by gain share (\texttt{feature\_importance.png})\\
  and the share of hours above \(\tau\) by clock hour (\texttt{breach\_by\_hour.png}).}
  \caption{Drivers of the forecast and the clock hours that cross the capacity threshold.}
  \label{fig:importance}
\end{figure}

The breach-by-hour profile is the basis for the reserve rule. Complete the
following statement from \texttt{breach\_rate\_by\_hour\_pct}. The design
target of the study is the evening window from 18:00 to 21:00, which is
expected to concentrate threshold crossings and therefore to need reserve
armed one to two hours earlier. If the measured profile peaks in a
neighbouring window, state that window and keep the same lead time.

\begin{quote}
On the hourly series, the hours from [HH:00] to [HH:00] show the highest
rate of demand above \(\tau\) ([XX]\% of hours in the busiest clock hour).
Reserve for that window should be committed one to two hours in advance.
Temperature and hour of day are the leading exogenous contributors to the
peak. Their LightGBM gain shares are [XX]\% and [XX]\%.
\end{quote}
